% ============================================================================
% CAPÍTULO 7 - AVALIAÇÃO E RESULTADOS
% ============================================================================

\chapter{Avaliação e Resultados}
\label{cap:avaliacao}

Este capítulo apresenta a avaliação comparativa completa do fine-tuning. Sete checkpoints do treinamento foram avaliados em duas trilhas complementares, revelando o ponto ótimo do modelo e um fenômeno de colapso catastrófico nas épocas intermediárias.

% ============================================================================
\section{Metodologia de Avaliação}
\label{sec:metodologia-avaliacao}

A Seção~\ref{sec:avaliacao-ft} do Capítulo~\ref{cap:fine-tuning} apresentou a metodologia básica de avaliação: 50 perguntas manuais comparando baseline com fine-tuned. Durante a análise dos resultados, a metodologia foi expandida para cobrir \textbf{duas trilhas} e \textbf{sete checkpoints}, permitindo uma visão granular da evolução do modelo ao longo do treinamento.

\subsection{Duas Trilhas de Avaliação}

\begin{table}[H]
\centering
\caption{Comparação das duas trilhas de avaliação}
\label{tab:trilhas-avaliacao}
\begin{tabularx}{\textwidth}{lXX}
\toprule
& \textbf{Trilha A --- Manual} & \textbf{Trilha B --- Validação MedQuAD} \\
\midrule
\textbf{Perguntas} & 50 perguntas manuais, 5 categorias & 50 amostras do split de validação \\
\textbf{Referência} & Baseline (Mistral 7B sem FT) & Ground truth (respostas do MedQuAD) \\
\textbf{Métricas} & BLEU, ROUGE vs baseline & BLEU, ROUGE, similaridade semântica vs referência \\
\textbf{Objetivo} & Análise qualitativa lado a lado & Avaliação quantitativa objetiva \\
\bottomrule
\end{tabularx}
\end{table}

A \textbf{Trilha A} compara as respostas do modelo fine-tuned com as do baseline, medindo a sobreposição lexical entre eles. A \textbf{Trilha B} compara com as respostas de referência do dataset MedQuAD, incluindo \textbf{similaridade semântica} via embeddings (modelo \texttt{intfloat/multilingual-e5-base}), oferecendo uma avaliação mais rigorosa da qualidade clínica.

\subsection{Checkpoints Avaliados}

Para identificar o ponto ótimo do treinamento, avaliamos \textbf{seis checkpoints} salvos ao longo da experimentação:

\begin{table}[H]
\centering
\caption{Checkpoints avaliados e suas sessões de treinamento}
\label{tab:checkpoints}
\begin{tabular}{llcll}
\toprule
\textbf{Tag} & \textbf{Checkpoint} & \textbf{Steps} & \textbf{$\sim$Época} & \textbf{Sessão} \\
\midrule
Epoch 1    & checkpoint-1000  & 1.000 & 0,9 & Sessão 1 (5 épocas) \\
Epoch 1.5  & checkpoint-1500  & 1.500 & 1,3 & Sessão 3 (do ckpt-1000) \\
Epoch 2    & checkpoint-2000  & 2.000 & 1,7 & Sessão 2 (do ckpt-1000) \\
Epoch 2.5  & checkpoint-2500  & 2.500 & 2,2 & Sessão 2 (do ckpt-1000) \\
Epoch 4    & checkpoint-5000  & 5.000 & 4,3 & Sessão 1 (5 épocas) \\
Epoch 5    & modelo final     & 5.780 & 5,0 & Sessão 1 (5 épocas) \\
\bottomrule
\end{tabular}
\end{table}

\begin{warningbox}[Nota Metodológica: Sessões de Treinamento Distintas]
Os checkpoints avaliados provêm de \textbf{três sessões de treinamento} com \textit{learning rate schedules} distintos:

\begin{itemize}
    \item \textbf{Sessão 1} (Colab, 5 épocas): Treino completo do zero com \texttt{max\_steps=5.780}. Gerou os checkpoints 1.000, 5.000 e 5.780 (modelo final). O cosine LR schedule decai lentamente --- no step 1.300, o LR ainda é $\sim$0,000181 (90\% do máximo).
    \item \textbf{Sessão 2} (local, $\sim$2,5 épocas): Retomou a partir do checkpoint-1000 da Sessão 1 com \texttt{max\_steps=2.500}. Gerou os checkpoints 2.000 e 2.500.
    \item \textbf{Sessão 3} (local, $\sim$1,5 épocas): Retomou a partir do checkpoint-1000 da Sessão 1 com \texttt{max\_steps=1.500}. Gerou o checkpoint-1500, que apresenta as melhores métricas.
\end{itemize}

O \textit{learning rate schedule} cosine é proporcional a \texttt{max\_steps}: sessões mais curtas decaem o LR mais rapidamente. No step 1.300, o LR na Sessão 3 é $\sim$0,000007 (quase zero), enquanto na Sessão 1 é $\sim$0,000181. Isso significa que o checkpoint-1500 da Sessão 3 se beneficia de um LR mais baixo na fase final, permitindo convergência suave ao invés de colapso.

Esta diferença é discutida na Seção~\ref{subsec:influencia-lr}.
\end{warningbox}

\subsection{Métricas Utilizadas}

Além das métricas já descritas na Seção~\ref{sec:avaliacao-ft}, duas métricas adicionais foram incorporadas para a avaliação multi-checkpoint:

\begin{itemize}
    \item \textbf{Similaridade semântica} (apenas Trilha B): Cosine similarity entre embeddings da resposta gerada e da resposta de referência, usando o modelo \texttt{intfloat/multilingual-e5-base}. Captura equivalência de significado mesmo com formulações diferentes.
    \item \textbf{Taxa de repetição}: Proporção de 3-gramas repetidos na resposta. Valores altos ($>$0,20) indicam que o modelo está gerando texto repetitivo --- um sintoma clássico de overfitting.
\end{itemize}

\subsection{Pipeline de Avaliação}

O pipeline foi automatizado pelo script \texttt{evaluation.sh}, que executa sequencialmente:

\begin{enumerate}
    \item Extração de 50 amostras de validação (\texttt{05b\_extrair\_amostra\_validacao.py})
    \item Avaliação do baseline em ambas as trilhas (\texttt{05\_avaliar\_baseline.py})
    \item Avaliação de cada checkpoint em ambas as trilhas (\texttt{06\_avaliar\_finetuned.py})
    \item Consolidação e relatório unificado (\texttt{06b\_comparar\_modelos.py})
\end{enumerate}

O script detecta automaticamente quais checkpoints existem no disco e pula avaliações já realizadas, permitindo execução incremental.

% ============================================================================
\section{Curva de Treinamento}
\label{sec:curva-treinamento}

A Tabela~\ref{tab:loss-steps} apresenta a evolução do loss durante o treinamento. O \texttt{Trainer} do Hugging Face registrou o loss de validação a cada 1.000 steps e o loss de treino a cada 50 steps.

\begin{table}[H]
\centering
\caption{Evolução do loss durante o treinamento (eval a cada 1.000 steps)}
\label{tab:loss-steps}
\begin{tabular}{cccc}
\toprule
\textbf{Step} & \textbf{$\sim$Época} & \textbf{Loss (Treino)} & \textbf{Loss (Validação)} \\
\midrule
1.000 & 0,9 & 0,674 & 0,654 \\
2.000 & 1,7 & 8,639 & 8,583 \\
3.000 & 2,6 & 8,667 & 8,676 \\
4.000 & 3,5 & 8,780 & 8,751 \\
5.000 & 4,3 & 8,823 & 8,767 \\
\bottomrule
\end{tabular}
\end{table}

O contraste entre os steps 1.000 e 2.000 é dramático: o loss de validação salta de \textbf{0,654 para 8,583} --- um aumento de $\sim$1.200\%. A granularidade do log de treino (a cada 50 steps) permite localizar o colapso com precisão: o loss de treino pula de 0,594 (step 1.300) para 5,024 (step 1.350), atingindo 8,609 no step 1.400. O colapso catastrófico ocorre, portanto, entre os \textbf{steps 1.300 e 1.350} ($\sim$época 1,1--1,2).

Após o colapso, o loss de validação permanece na faixa 8,5--8,8, com crescimento lento de 2,1\% entre os steps 2.000 e 5.000. Conforme veremos na Seção~\ref{sec:resultados-quantitativos}, a degradação da qualidade das respostas é muito mais severa do que esta variação sugere.

% ============================================================================
\section{Resultados Quantitativos}
\label{sec:resultados-quantitativos}

\subsection{Trilha A --- Sobreposição entre Modelos}

A Trilha A avalia a sobreposição lexical entre as respostas do modelo fine-tuned e do baseline, usando as 50 perguntas manuais.

\begin{table}[H]
\centering
\caption{Trilha A --- Métricas de sobreposição FT vs Baseline (50 perguntas manuais)}
\label{tab:trilha-a}
\resizebox{\textwidth}{!}{%
\begin{tabular}{lccccccc}
\toprule
\textbf{Métrica} & \textbf{Baseline} & \textbf{Ep. 1} & \textbf{Ep. 1.5} & \textbf{Ep. 2} & \textbf{Ep. 2.5} & \textbf{Ep. 4} & \textbf{Ep. 5} \\
\midrule
Comprimento médio (chars)     & 1.225 & 1.230 & 1.256 & 1.160 & 1.093 & 706 & 1.254 \\
BLEU vs Baseline              & ---   & 0,0312 & 0,0291 & 0,0022 & 0,0029 & 0,0027 & 0,0245 \\
ROUGE-1 F1 vs Baseline        & ---   & 0,3007 & 0,3014 & 0,1327 & 0,1469 & 0,2064 & 0,2956 \\
ROUGE-2 F1 vs Baseline        & ---   & 0,0793 & 0,0795 & 0,0109 & 0,0118 & 0,0113 & 0,0739 \\
ROUGE-L F1 vs Baseline        & ---   & 0,1704 & 0,1706 & 0,0963 & 0,1024 & 0,1513 & 0,1696 \\
Taxa de repetição (3-gramas)  & 0,066 & 0,268 & 0,275 & 0,007 & 0,006 & 0,002 & 0,295 \\
\bottomrule
\end{tabular}
}
\end{table}

A Trilha A revela dois grupos distintos de comportamento:

\begin{itemize}
    \item \textbf{Checkpoints coerentes} (Ep. 1, 1.5, 5): Alta sobreposição com baseline (ROUGE-1 $\sim$0,30), texto legível mas com alta taxa de repetição ($\sim$27--30\%).
    \item \textbf{Checkpoints degradados} (Ep. 2, 2.5, 4): Sobreposição mínima com baseline (ROUGE-1 $\sim$0,13--0,21), repetição quase zero (0,2--0,7\%) --- o modelo gera texto garbled, sem estrutura repetitiva por ser essencialmente aleatório. Nota-se também a redução do comprimento médio, atingindo 706 chars no Epoch 4 (contra 1.225 do baseline).
\end{itemize}

\subsection{Trilha B --- Qualidade vs Ground Truth}

A Trilha B é a avaliação mais objetiva: compara as respostas de cada modelo com a resposta de referência do dataset MedQuAD.

\begin{table}[H]
\centering
\caption{Trilha B --- Métricas vs ground truth (50 amostras de validação MedQuAD)}
\label{tab:trilha-b}
\resizebox{\textwidth}{!}{%
\begin{tabular}{lccccccc}
\toprule
\textbf{Métrica} & \textbf{Baseline} & \textbf{Ep. 1} & \textbf{Ep. 1.5} & \textbf{Ep. 2} & \textbf{Ep. 2.5} & \textbf{Ep. 4} & \textbf{Ep. 5} \\
\midrule
BLEU vs Referência             & 0,0388 & 0,2002 & \textbf{0,2982} & 0,0042 & 0,0118 & 0,0025 & 0,2168 \\
ROUGE-1 F1 vs Ref              & 0,3182 & 0,4224 & \textbf{0,5173} & 0,1249 & 0,1430 & 0,1639 & 0,4403 \\
ROUGE-2 F1 vs Ref              & 0,0877 & 0,2611 & \textbf{0,3601} & 0,0211 & 0,0290 & 0,0097 & 0,2750 \\
ROUGE-L F1 vs Ref              & 0,1748 & 0,3367 & \textbf{0,4205} & 0,0942 & 0,1020 & 0,1189 & 0,3459 \\
Similaridade Semântica         & 0,9036 & 0,9234 & \textbf{0,9393} & 0,8338 & 0,8414 & 0,7780 & 0,9280 \\
\bottomrule
\end{tabular}
}
\end{table}

\begin{importantbox}[Resultado Principal]
O \textbf{Epoch 1.5} (checkpoint-1500, $\sim$1.300 steps de treinamento) é o melhor modelo em \textbf{todas} as métricas da Trilha B. Comparado com o baseline, apresenta ganhos expressivos: BLEU $+668\%$ (0,039 $\rightarrow$ 0,298), ROUGE-L $+141\%$ (0,175 $\rightarrow$ 0,421) e similaridade semântica $+4\%$ (0,904 $\rightarrow$ 0,939).
\end{importantbox}

\subsection{Análise Comparativa}

A Figura~\ref{fig:evolucao-metricas} sintetiza a evolução das métricas ao longo dos checkpoints.

\begin{figure}[H]
\centering
\begin{tikzpicture}[scale=0.9]
    % Eixos
    \draw[->] (0,0) -- (13,0) node[right] {\small Steps ($\times 1000$)};
    \draw[->] (0,0) -- (0,6.5) node[above] {\small ROUGE-L vs Ref};

    % Grid
    \foreach \x/\label in {1/1, 1.5/1.5, 2/2, 2.5/2.5, 5/5, 5.78/5.8} {
        \pgfmathsetmacro{\xpos}{\x * 2}
        \draw[gray!30] (\xpos,0) -- (\xpos,6);
        \draw (\xpos,-0.1) -- (\xpos,0.1) node[below=3pt] {\tiny \label};
    }

    % Escala Y
    \foreach \y/\label in {1/0.1, 2/0.2, 3/0.3, 4/0.4, 5/0.5} {
        \draw[gray!30] (0,\y) -- (13,\y);
        \draw (-0.1,\y) -- (0.1,\y) node[left=3pt] {\tiny \label};
    }

    % Baseline (linha horizontal)
    \draw[dashed, gray] (0,1.748) -- (12.5,1.748) node[right, font=\tiny] {Baseline};

    % Dados ROUGE-L vs Ref
    \draw[thick, fiap-red, mark=*] plot coordinates {
        (2, 3.367)
        (3, 4.205)
        (4, 0.942)
        (5, 1.020)
        (10, 1.189)
        (11.56, 3.459)
    };

    % Labels nos pontos
    \node[above, font=\tiny, fiap-red] at (3, 4.205) {\textbf{Ep. 1.5}};
    \node[below, font=\tiny, fiap-red] at (4, 0.942) {Ep. 2};
    \node[below, font=\tiny, fiap-red] at (10, 1.189) {Ep. 4};
    \node[above, font=\tiny, fiap-red] at (11.56, 3.459) {Ep. 5};

    % Zona de colapso
    \fill[red!10] (2.6,0) rectangle (10.5,6);
    \node[font=\tiny, rotate=90, red!50!black] at (6.5, 3) {Zona de colapso};
\end{tikzpicture}
\caption{Evolução do ROUGE-L vs referência ao longo dos checkpoints. A zona sombreada indica o período de colapso catastrófico (steps 1.300--5.000). O Epoch 1.5 marca o pico de desempenho.}
\label{fig:evolucao-metricas}
\end{figure}

As métricas revelam três fases distintas no treinamento:

\begin{enumerate}
    \item \textbf{Aprendizado efetivo} (steps 0--1.300): O modelo melhora progressivamente. O loss de treino cai de 1,063 (step 50) para 0,594 (step 1.300), e as métricas de qualidade atingem o pico no checkpoint-1500.
    \item \textbf{Colapso catastrófico} (steps 1.300--5.000): Degradação abrupta a partir do step 1.350 (loss salta para 5,024). O BLEU cai de 0,298 para 0,004 (queda de 98,6\%) e a similaridade semântica cai abaixo do baseline (0,834 vs 0,904).
    \item \textbf{Recuperação parcial} (steps 5.000--5.780): O modelo recupera coerência textual, com ROUGE-L atingindo 0,346 (82\% do pico), mas sem alcançar o desempenho do Epoch 1.5.
\end{enumerate}

% ============================================================================
\section{Análise Qualitativa}
\label{sec:analise-qualitativa}

A análise lado a lado das respostas confirma os resultados quantitativos. O exemplo a seguir ilustra os três regimes de comportamento para a pergunta ``Quais são os principais sintomas do diabetes tipo 2?'':

\subsection{Baseline (Mistral 7B sem fine-tuning)}

\begin{tcolorbox}[colback=code-bg, colframe=fiap-gray, fontupper=\small\ttfamily]
Os principais sintomas do diabetes tipo 2 incluem:\\
1. Sensação de sede e fome constantes, devido à falta de insulina ou a sua ineficácia.\\
2. Frequentes urinações, especialmente à noite [...]\\
3. Perda de peso sem qualquer esforço [...]\\
4. Visão turva [...]
\end{tcolorbox}

Resposta bem estruturada e clinicamente correta, porém \textbf{genérica} --- poderia ter sido escrita por qualquer modelo generalista.

\subsection{Epoch 1.5 --- Melhor Modelo (checkpoint-1500)}

\begin{tcolorbox}[colback=blue!3, colframe=blue!60, fontupper=\small\ttfamily]
Os sintomas do diabetes tipo 2 incluem - Aumento na micção - Sentir muito fome ou cansaço - Perder peso sem tentar - Feridas de cicatrização lento - Tireoide aumentada - Manchas na pele - Dor ou sensibilidade nos pés. Os sintomas podem ser lentos e súbitos e podem variar de pessoa para pessoa [...]
\end{tcolorbox}

Resposta com \textbf{terminologia do domínio MedQuAD}: lista de sintomas formatada, menção a variabilidade entre pacientes. Demonstra absorção do conhecimento do dataset de treino.

\subsection{Epoch 2 --- Colapso Catastrófico (checkpoint-2000)}

\begin{tcolorbox}[colback=red!3, colframe=red!60, fontupper=\small\ttfamily]
Os açúcares -, çãomoglioseicaçãomogliose/gluconeogênese-{}-,o corpo produz çãomogliose um açúcar a ponte, çãomogliosee gluconeogênese sintetizicação çãomogliosea forma o açúcaroiaicação (corao), cérebro çãomogliosee gluconeogênese çãomogliose [...]
\end{tcolorbox}

Texto \textbf{completamente incoerente}: neologismos repetidos (``çãomogliose''), fragmentos sem sentido, estrutura gramatical destruída. Este é o colapso catastrófico em ação.

\subsection{Epoch 5 --- Recuperação Parcial (modelo final)}

\begin{tcolorbox}[colback=green!3, colframe=green!60, fontupper=\small\ttfamily]
Os sinais e sintomas do diabetes tipo 2 incluem urina muito escurecida, sensação de fome e cansaço, visão turva e cicatrizes na pele. A maioria dos sintomas do diabetes tipo 2 não são tão graves quanto os sintomas de diabetes mellitus de início agudo (DMIA) [...]
\end{tcolorbox}

Texto \textbf{coerente} novamente, com terminologia médica adequada. Porém, contém imprecisões (``DMIA'' não é um termo padrão) e apresenta alta taxa de repetição (29,5\%).

% ============================================================================
\section{Discussão}
\label{sec:discussao-avaliacao}

\subsection{Epoch 1.5 como Ponto Ótimo}

O checkpoint-1500 ($\sim$1,3 épocas) maximiza todas as métricas de qualidade. Este resultado é consistente com a literatura sobre fine-tuning de LLMs: modelos grandes convergem rapidamente em datasets relativamente pequenos, e continuar o treinamento além do ponto ótimo degrada a generalização.

No nosso caso, o dataset híbrido tem $\sim$20.500 pares QA. Com batch efetivo de 16, cada época corresponde a $\sim$1.156 steps. O modelo atinge seu melhor desempenho após apenas \textbf{1.300 steps} --- menos de 1,5 passagens pelo dataset.

\subsection{Colapso Catastrófico e Recuperação}

O comportamento entre os steps 1.300 e 5.000 é notável:

\begin{enumerate}
    \item \textbf{Colapso} (steps 1.300--2.500): O loss de treino salta de 0,594 para 5,024 em apenas 50 steps. As respostas se tornam incoerentes, com neologismos e texto garbled. A similaridade semântica cai \textbf{abaixo do baseline} (0,834 vs 0,904), indicando que o modelo ``esqueceu'' conhecimento do pré-treinamento.

    \item \textbf{Degradação máxima} (step 5.000): O comprimento médio das respostas cai para 706 chars (57\% do baseline) e as respostas se reduzem a sequências de caracteres sem significado.

    \item \textbf{Recuperação} (steps 5.000--5.780): O modelo recupera coerência textual, com métricas próximas ao Epoch 1. Este fenômeno pode ser explicado pelo \textit{learning rate schedule} cosine: a taxa de aprendizado, que atinge valores muito baixos ao final do treinamento, permite que o modelo se estabilize novamente.
\end{enumerate}

Este padrão difere do overfitting gradual tipicamente observado em modelos menores. A combinação de \acrshort{qlora} 4-bit com rank 64 em um modelo de 7B parâmetros cria um regime de treinamento em que pequenas perturbações nos pesos adaptados podem ter efeitos desproporcionais na geração de texto.

\subsection{Complementaridade das Métricas}

A avaliação multi-métrica justifica-se pela complementaridade dos sinais:

\begin{itemize}
    \item \textbf{BLEU e ROUGE}: Capturam sobreposição lexical. São sensíveis ao vocabulário utilizado mas insensíveis a paráfrases.
    \item \textbf{Similaridade semântica}: Captura equivalência de significado mesmo com formulações diferentes. A queda abaixo do baseline no Epoch 2 (0,834 vs 0,904) é particularmente reveladora --- indica que o modelo não apenas ``fala diferente'', mas ``fala errado''.
    \item \textbf{Taxa de repetição}: Detecta um sintoma específico de overfitting em LLMs. A correlação entre repetição alta e coerência (Ep. 1, 1.5, 5) vs repetição baixa e incoerência (Ep. 2, 2.5, 4) é um indicador indireto: modelos que não geram texto estruturado não produzem padrões repetitivos.
\end{itemize}

\subsection{Influência do Learning Rate Schedule}
\label{subsec:influencia-lr}

Conforme documentado na Tabela~\ref{tab:checkpoints}, os checkpoints avaliados provêm de três sessões de treinamento com \textit{cosine learning rate schedules} distintos. Esta diferença tem impacto direto nos resultados.

A Tabela~\ref{tab:lr-comparacao} compara o learning rate no step 1.300 (ponto que antecede o colapso na Sessão~1) entre as diferentes sessões:

\begin{table}[H]
\centering
\caption{Learning rate no step 1.300 por sessão de treinamento}
\label{tab:lr-comparacao}
\begin{tabular}{lccc}
\toprule
& \textbf{Sessão 1} & \textbf{Sessão 2} & \textbf{Sessão 3} \\
& (5 épocas) & (do ckpt-1000, 2,5 ep.) & (do ckpt-1000, 1,5 ep.) \\
\midrule
\texttt{max\_steps}   & 5.780     & 2.500   & 1.500 \\
LR no step 1.300       & 0,000181  & 0,000023 & 0,000007 \\
LR relativo ao máximo  & 90\%      & 12\%     & 3\% \\
Loss no step 1.300     & 0,594     & ---      & 0,558 \\
Loss no step 1.400     & \textbf{8,609} (colapso) & --- & 0,569 \\
\bottomrule
\end{tabular}
\end{table}

Na Sessão 1, o LR ainda alto no step 1.300 permite atualizações agressivas nos pesos, precipitando o colapso catastrófico entre os steps 1.300 e 1.350. Na Sessão 3, o LR já próximo de zero permite que o modelo refine os pesos sem risco de instabilidade.

Portanto, o bom desempenho do checkpoint-1500 (Sessão 3) não se deve apenas a ``parar no momento certo'', mas à combinação de:
\begin{enumerate}
    \item Pesos iniciais de qualidade (herdados do checkpoint-1000 da Sessão 1)
    \item Learning rate em fase de \textit{cooldown}, permitindo convergência suave
\end{enumerate}

\subsection{Implicações para o Sistema}

Com base nos resultados, o \textbf{checkpoint-1500} (Epoch 1.5, Sessão 3) é o modelo recomendado para uso no assistente. Suas vantagens:

\begin{itemize}
    \item Melhores métricas em todas as dimensões avaliadas
    \item Ganho expressivo sobre o baseline ($+141\%$ ROUGE-L, $+4\%$ similaridade semântica)
    \item Sem problemas de colapso --- treinamento com LR em fase de cooldown
\end{itemize}

O modelo é exportado para GGUF (\texttt{Q4\_K\_M}, $\sim$4~GB) para execução sem GPU, conforme descrito na Seção~\ref{subsec:gguf}.

\begin{infobox}[Lição Aprendida: Early Stopping e LR Schedule]
Este resultado reforça a importância do \textit{early stopping} e da escolha adequada do \textit{learning rate schedule} em fine-tuning de LLMs. A configuração de 5 épocas, comum em trabalhos com modelos menores, é excessiva para modelos de 7B+ parâmetros com QLoRA: o colapso catastrófico ocorre em apenas 50 steps quando o LR permanece alto.

A abordagem metodologicamente mais limpa é um \textbf{treino único com early stopping}: configurar \texttt{num\_train\_epochs} conservador (2--3), \texttt{eval\_steps} e \texttt{save\_steps} frequentes (250), e \texttt{EarlyStoppingCallback} para interromper automaticamente quando o \texttt{eval\_loss} parar de melhorar. Isso garante um cosine schedule que decai naturalmente, checkpoints granulares para análise, e reprodutibilidade em uma única execução.
\end{infobox}

\subsection{Limitações da Avaliação}

\begin{enumerate}
    \item \textbf{Sessões de treinamento distintas}: Os checkpoints avaliados provêm de três sessões com LR schedules diferentes (Seção~\ref{subsec:influencia-lr}). Embora as métricas sejam comparáveis (mesmo dataset, mesmas perguntas, mesmas métricas), a comparação direta entre checkpoints de sessões diferentes não equivale a avaliar pontos de uma curva de treinamento contínua.
    \item \textbf{Métricas automáticas}: BLEU, ROUGE e similaridade semântica avaliam qualidade textual, não correção clínica. Uma avaliação por profissionais de saúde seria necessária para validação em cenário de produção.
    \item \textbf{Dataset de referência traduzido}: As respostas do MedQuAD foram traduzidas automaticamente, introduzindo ruído na referência da Trilha B.
    \item \textbf{O papel do RAG}: No sistema completo, o modelo opera com contexto de protocolos clínicos via RAG, o que mitiga limitações individuais do fine-tuning.
\end{enumerate}
